% Source: physical pp. 198--201 (printed pp. 175--178). Coverage:
% Problems 1--5 and References 1--21, ending before Chapter 28.
% Figures/tables/footnotes/dividers: none.

\chapterbackmatter{Problems}

\begin{enumerate}
  \item To second order in gauge coupling constants, calculate the
  components of the superfields $\Omega^A$ that are needed to put the
  gauge superfield $V^A$ in Wess--Zumino gauge by the transformation
  (27.1.12).

  \item Show that the most general chiral spinor superfield $W_\alpha$
  satisfying the condition (27.2.20) has components $f_{\mu\nu}$
  satisfying the homogeneous Maxwell equations
  $\epsilon^{\mu\nu\rho\sigma}\partial_\rho f_{\mu\nu}=0$. What
  conditions on the other components of $W_\alpha$ are imposed by
  Eq. (27.2.20)?

  \item Consider a general renormalizable $N=1$ supersymmetric gauge
  theory with an $SU(2)$ gauge group and a single chiral superfield
  belonging to the 3-vector representation of $SU(2)$. What is the most
  general superpotential for this theory? Construct the Lagrangian
  density of the whole theory explicitly. Eliminate the auxiliary
  fields. Show that supersymmetry is not broken in this theory. What are
  the masses of the particles of this theory?

  \item Express the gaugino and chiral fermion fields in the
  supersymmetric version of quantum electrodynamics described in
  Section 27.5 in terms of the goldstino field and other spinor fields
  of definite mass.

  \item Consider the renormalizable $N=2$ supersymmetric theory with an
  $SU(3)$ gauge symmetry and no hypermultiplets. What are the values of
  the scalar fields for which the potential vanishes? What are the
  massless gauge fields for non-zero values of these scalars? Calculate
  the central charge in terms of the quantities to which these massless
  gauge fields are coupled.
\end{enumerate}

\chapterbackmatter{References}

\begin{enumerate}
  \item \phantomsection\label{ch27-ref-1}
  Supersymmetry was applied first to Abelian gauge theories without
  using the superfield formalism by J.~Wess and B.~Zumino,
  \textit{Nucl. Phys.} B\textbf{78}, 1 (1974). It was then extended to
  non-Abelian gauge theories by S.~Ferrara and B.~Zumino,
  \textit{Nucl. Phys.} B\textbf{79}, 413 (1974); A.~Salam and
  J.~Strathdee, \textit{Phys. Lett.} \textbf{51B}, 353 (1974). These
  references are reprinted in \textit{Supersymmetry}, S.~Ferrara, ed.
  (North Holland/World Scientific, Amsterdam/Singapore, 1987).

  \item \phantomsection\label{ch27-ref-2}
  P.~Fayet and J.~Iliopoulos, \textit{Phys. Lett.} \textbf{51B}, 461
  (1974). This reference is reprinted in \textit{Supersymmetry},
  Ref.~\hyperref[ch27-ref-1]{1}.

  \item \phantomsection\label{ch27-ref-3}
  S.~Weinberg, \textit{Phys. Rev. Lett.} \textbf{80}, 3702 (1998).

  \item \phantomsection\label{ch27-ref-4}
  S.~Ferrara, L.~Girardello, and F.~Palumbo, \textit{Phys. Rev.}
  D\textbf{20}, 403 (1979). This article is reprinted in
  \textit{Supersymmetry}, Ref.~\hyperref[ch27-ref-1]{1}. Special cases
  of this sum rule had been given by P.~Fayet, \textit{Phys. Lett.}
  \textbf{84B}, 416 (1979).

  \item \phantomsection\label{ch27-ref-5}
  M.~T.~Grisaru, W.~Siegel, and M.~Ro\v{c}ek, \textit{Nucl. Phys.}
  B\textbf{159}, 429 (1979).

  \item \phantomsection\label{ch27-ref-6}
  N.~Seiberg, \textit{Phys. Lett.} B\textbf{318}, 469 (1993).

  \item \phantomsection\label{ch27-ref-7}
  This was first proved in the supergraph formalism by W.~Fischler,
  H.~P.~Nilles, J.~Polchinski, S.~Raby, and L.~Susskind,
  \textit{Phys. Rev. Lett.} \textbf{47}, 757 (1981). The proof presented
  here was given by M.~Dine, in \textit{Fields, Strings, and Duality:
  TASI 96}, C.~Efthimiou and B.~Greene, eds. (World Scientific,
  Singapore, 1997); S.~Weinberg, Ref.~\hyperref[ch27-ref-3]{3}.

  \item \phantomsection\label{ch27-ref-8}
  A detailed proof of the statement that superrenormalizable terms that
  violate some global symmetry do not introduce infinite symmetry
  breaking radiative corrections to the coefficients of renormalizable
  interactions is given by K.~Symanzik, in \textit{Carg\`ese Lectures
  in Physics}, Vol.~5, D.~Bessis, ed. (Gordon and Breach, New York,
  1972). This is discussed briefly in the footnote on p.~507 of Volume
  I.

  \item \phantomsection\label{ch27-ref-9}
  L.~Girardello and M.~T.~Grisaru, \textit{Nucl. Phys.} B\textbf{194},
  65 (1982), reprinted in \textit{Supersymmetry},
  Ref.~\hyperref[ch27-ref-1]{1}; K.~Harada and N.~Sakai,
  \textit{Prob. Theor. Phys.} \textbf{67}, 67 (1982).

  \item \phantomsection\label{ch27-ref-10}
  For a review, see S.~Weinberg, \textit{Rev. Mod. Phys.} \textbf{61},
  1--23 (1989).

  \item \phantomsection\label{ch27-ref-11}
  B.~de Wit and D.~Z.~Freedman, \textit{Phys. Rev.} D\textbf{12}, 2286
  (1975). This article is reprinted in \textit{Supersymmetry},
  Ref.~\hyperref[ch27-ref-1]{1}.

  \item \phantomsection\label{ch27-ref-12}
  The first examples of gauge theories with $N=2$ extended
  supersymmetry were given by P.~Fayet, \textit{Nucl. Phys.}
  B\textbf{113}, 135 (1976); reprinted in \textit{Supersymmetry},
  Ref.~\hyperref[ch27-ref-1]{1}. The approach presented here is similar
  to that of Fayet. A superfield formalism was subsequently given by
  R.~Grimm, M.~Sohnius, and J.~Scherk, \textit{Nucl. Phys.}
  B\textbf{113}, 77 (1977). Gauge theories with $N=2$ and $N=4$
  supersymmetry in four spacetime dimensions were constructed by
  dimensional reduction of higher-dimensional theories with simple
  supersymmetry by L.~Brink, J.~H.~Schwarz, and J.~Scherk,
  \textit{Nucl. Phys.} B\textbf{113}, 77 (1977); M.~F.~Sohnius,
  K.~S.~Stelle, and P.~C.~West, \textit{Nucl. Phys.} B\textbf{113}, 127
  (1980). For other approaches, see M.~F.~Sohnius, \textit{Nucl. Phys.}
  B\textbf{138}, 109 (1979); A.~Halperin, E.~A.~Ivanov, and
  V.~I.~Ogievetsky, \textit{Prima JETP} \textbf{33}, 176 (1981);
  P.~Breitenlohner and M.~F.~Sohnius, \textit{Nucl. Phys.}
  B\textbf{178}, 151 (1981); P.~Howe, K.~S.~Stelle, and
  P.~K.~Townsend, \textit{Nucl. Phys.} B\textbf{214}, 519 (1983).

  \item \phantomsection\label{ch27-ref-13}
  E.~Witten and D.~Olive, \textit{Phys. Lett.} \textbf{78B}, 97 (1978).
  Also see H.~Osborn, \textit{Phys. Lett.} \textbf{83B}, 321 (1979).

  \item \phantomsection\label{ch27-ref-14}
  E.~B.~Bogomol'nyi, \textit{Sov. J. Nucl. Phys.} \textbf{24}, 449
  (1976).

  \item \phantomsection\label{ch27-ref-15}
  D.~Zwanziger, \textit{Phys. Rev.} \textbf{176}, 1480, 1489 (1968);
  J.~Schwinger, \textit{Phys. Rev.} \textbf{144}, 1087 (1966);
  \textbf{173}, 1536 (1968); B.~Julia and A.~Zee,
  \textit{Phys. Rev.} D\textbf{11}, 2227 (1974); F.~A.~Bais and
  J.~R.~Primack, \textit{Phys. Rev.} D\textbf{13}, 819 (1975). (The
  last article was incorrectly attributed to Julia and Zee in Chapter
  23 of the first printings of Volume II.)

  \item \phantomsection\label{ch27-ref-16}
  M.~K.~Prasad and C.~M.~Sommerfield, \textit{Phys. Rev. Lett.}
  \textbf{35}, 760 (1975); E.~B.~Bogomol'nyi,
  Ref.~\hyperref[ch27-ref-14]{14}; S.~Coleman, S.~Parke, A.~Neveu, and
  C.~M.~Sommerfield, \textit{Phys. Rev.} D\textbf{15}, 544 (1977).

  \item \phantomsection\label{ch27-ref-17}
  This was noted by C.~Montonen and D.~Olive, \textit{Phys. Lett.}
  \textbf{72B}, 117 (1977). The absence of one-loop corrections to the
  masses was shown by A.~D'Adda, R.~Horsley, and P.~Di Vecchia,
  \textit{Phys. Lett.} \textbf{76B}, 298 (1978).

  \item \phantomsection\label{ch27-ref-18}
  Obstacles to formulating theories with auxiliary fields for $N=4$
  supersymmetry have been analyzed by W.~Siegel and M.~Ro\v{c}ek,
  \textit{Phys. Lett.} \textbf{105B}, 275 (1981).

  \item \phantomsection\label{ch27-ref-19}
  The finiteness of the $N=4$ theory was shown by M.~F.~Sohnius and
  P.~C.~West, \textit{Nucl. Phys.} B\textbf{100}, 245 (1981);
  P.~S.~Howe, K.~S.~Stelle, and P.~Townsend, \textit{Nucl. Phys.}
  B\textbf{214}, 519 (1983); S.~Mandelstam, \textit{Nucl. Phys.}
  B\textbf{213}, 149 (1983); L.~Brink, O.~Lindgren, and
  B.~E.~W.~Nilsson, \textit{Nucl. Phys.} B\textbf{212}, 401 (1983);
  \textit{Phys. Lett.} \textbf{123B}, 328 (1983). The proof of
  finiteness was extended to non-perturbative effects by N.~Seiberg,
  \textit{Phys. Lett.} B\textbf{206}, 75 (1988). Also see S.~Kovacs,
  hep-th/9902047, to be published. There is a wide class of
  ultraviolet-finite $N=2$ theories; see P.~S.~Howe, K.~S.~Stelle, and
  P.~C.~West, \textit{Phys. Lett.} B\textbf{124}, 55 (1983).

  \item \phantomsection\label{ch27-ref-20}
  H.~Osborn, Ref.~\hyperref[ch27-ref-13]{13}.

  \item \phantomsection\label{ch27-ref-21}
  A.~Sen, \textit{Phys. Lett.} B\textbf{329}, 217 (1994); C.~Vafa and
  E.~Witten, \textit{Nucl. Phys.} B\textbf{431}, 3 (1994);
  L.~Girardello, A.~Giveon, M.~Porrati, and A.~Zaffaroni,
  \textit{Phys. Lett.} B\textbf{334}, 331 (1994).
\end{enumerate}
